% Rebuttal to Reviewer Y2rf
% ACM-BCB 2026 Submission: Structure-Aware Prediction of PROTAC-Mediated
% Protein Degradability via Graph Neural Networks
%
% Compile: pdflatex rebuttal_Y2rf.tex && pdflatex rebuttal_Y2rf.tex

\documentclass[11pt]{article}
\usepackage[margin=1in]{geometry}
\usepackage{xcolor}
\usepackage{booktabs}
\usepackage{enumitem}
\usepackage{hyperref}
\usepackage{amsmath}

% Styling
\definecolor{revblue}{RGB}{0,70,170}
\newcommand{\reviewerbox}[2]{%
  \noindent\fcolorbox{black}{gray!10}{\parbox{\dimexpr\linewidth-2\fboxsep-2\fboxrule}{%
    \textbf{#1}: #2}}%
  \vspace{4pt}%
}
\newcommand{\resp}{\noindent\textbf{Response:}~}
\newcommand{\change}[1]{\textcolor{revblue}{[Paper reference: #1]}}

\setlength{\parskip}{6pt}
\setlength{\parindent}{0pt}

\begin{document}

\begin{center}
{\Large\bfseries Response to Reviewer Y2rf}\\[6pt]
{\large Structure-Aware Prediction of PROTAC-Mediated Protein Degradability\\via Graph Neural Networks}\\[6pt]
ACM-BCB 2026 Submission\\[4pt]
\textit{Rating: 2 --- Reject}
\end{center}

\bigskip

We appreciate the reviewer's detailed and rigorous critique. We have addressed every major and minor concern and believe the revised manuscript warrants reconsideration. All changes are \textcolor{revblue}{highlighted in blue} in the revised paper. Line numbers below refer to the revised \texttt{paper.tex}.

\bigskip
\hrule
\bigskip

%% ====================================================================
%%  M1
%% ====================================================================

\reviewerbox{M1}{The main result should be multi-seed average, not best seed. Rewrite abstract, results, and conclusion accordingly.}

\resp Done. Throughout the revised manuscript, the multi-seed average is the primary reported result and the best seed is secondary:

\begin{itemize}[topsep=2pt,itemsep=1pt]
  \item \textbf{Abstract} (line 65): Leads with ``0.646$\pm$0.124 AUROC on target-unseen evaluation (best seed: 0.7449)''---average first, best seed parenthetical.
  \change{Line 65.}

  \item \textbf{Abstract closing} (line 68): ``However, the high seed variance (std = 0.124) and limited E3 coverage require ensembling for reliable deployment.''
  \change{Line 68.}

  \item \textbf{Introduction} (lines 135--136): Reports both 3-seed (0.646$\pm$0.124) and 6-seed (0.603$\pm$0.097) means as primary numbers; notes statistical non-significance ($p = 0.556$) and VHL failure (0.396) immediately.
  \change{Lines 135--136.}

  \item \textbf{Table~1} (lines 441--458): Now contains both 3-seed and 6-seed rows, with ``Best Seed'' as a separate column rather than the headline number.
  \change{Lines 441--458 (Table~1, \texttt{tab:main}).}

  \item \textbf{Results text} (lines 460--462): ``The improved model achieves a 3-seed mean of 0.646\ldots and a more conservative 6-seed mean of 0.603\ldots The gradient boosting baseline achieves 0.607; the 3-seed advantage (+6.4\%) is not statistically significant (p = 0.259)\ldots The best seed (0.7449, +23\% over baseline) reflects favorable initialization and should not be expected reliably.''
  \change{Lines 460--462.}

  \item \textbf{Conclusion} (line 783): ``On target-unseen evaluation, the 3-seed mean is 0.646$\pm$0.124 (best: 0.7449) and the 6-seed mean is 0.603$\pm$0.097; neither is statistically superior to gradient boosting (0.607)\ldots''
  \change{Lines 783--784.}
\end{itemize}

\bigskip

%% ====================================================================
%%  M2
%% ====================================================================

\reviewerbox{M2}{Target-unseen performance is unstable. Report more seeds and provide confidence intervals. Discuss average performance rather than peak.}

\resp We trained 3 additional seeds (789, 1011, 1213), yielding AUROCs of 0.657, 0.520, and 0.502 respectively. The 6-seed results are:

\begin{center}
\begin{tabular}{lcc}
\toprule
\textbf{Estimate} & \textbf{AUROC} & \textbf{95\% CI} \\
\midrule
3-seed mean & 0.646 $\pm$ 0.124 & [0.506, 0.745] \\
6-seed mean & 0.603 $\pm$ 0.097 & [0.532, 0.682] \\
Gradient Boosting & 0.607 & --- \\
\bottomrule
\end{tabular}
\end{center}

The 6-seed mean (0.603) is marginally \emph{below} the gradient boosting baseline (0.607, bootstrap $p=0.556$). We report this honestly and prominently:

\begin{quote}\itshape
``This result is important to report honestly: the original 3-seed mean (0.646) was moderately favorable; the 6-seed analysis demonstrates that DegradoMap's advantage over hand-crafted feature baselines is not robust across all initializations.''
\end{quote}

We explicitly recommend ensembling across $\geq$6 seeds for deployment (line 624) and note that individual seeds span a 0.25-AUROC range (0.502--0.745).

\change{Lines 619--625 (``Six-seed extended validation'' paragraph in \S4.3); Table~1 lines 441--458 (6-seed row added); Figure~4 at line 632 (all 6 seeds visualized); Seed sensitivity paragraph lines 615--617.}

\bigskip

%% ====================================================================
%%  M3
%% ====================================================================

\reviewerbox{M3}{The lysine-related biological claim is not supported by ablation---removing the lysine indicator \emph{improves} performance. Explain or retract.}

\resp We now discuss this ``lysine indicator paradox'' directly in the Discussion (\S5) with concrete overfitting evidence:

\begin{itemize}[topsep=2pt,itemsep=1pt]
  \item With lysine indicator: val 0.584, test 0.477 (val-test gap: $-$0.106 $\to$ \textbf{overfitting}).
  \item Without lysine indicator: val 0.541, test 0.578 (val-test gap: $+$0.038 $\to$ generalizes).
\end{itemize}

This is a classic overfitting pattern: the lysine indicator improves validation but hurts test performance. We propose three explanations:

\begin{enumerate}[topsep=2pt,itemsep=1pt]
  \item The lysine-weighted \emph{pooling mechanism} (Eq.~3, line 247) already encodes lysine information architecturally, making the binary node feature redundant and causing overfitting to training-set lysine distributions.
  \item The model learns that not all lysines are equal---the binary indicator oversimplifies, while the pooling weights discriminate based on structural context (SASA, disorder, local geometry).
  \item The indicator may introduce a confounding bias: training targets have specific lysine distributions that don't generalize to unseen proteins.
\end{enumerate}

We conclude: ``architectural inductive biases (lysine pooling) are more effective than explicit feature annotation for this task,'' and note that current feature design is suboptimal.

\change{Lines 751--758 (``The lysine indicator paradox'' paragraph in \S5); Figure~7 at line 762 (lysine paradox visualization); Limitation item~7 at line 772 (``Lysine indicator counterproductive'').}

\bigskip

%% ====================================================================
%%  M4
%% ====================================================================

\reviewerbox{M4}{The E3 generalization claim should be narrower. Evidence supports mainly CRBN$\leftrightarrow$VHL transfer.}

\resp Agreed and revised throughout:

\begin{itemize}[topsep=2pt,itemsep=1pt]
  \item \textbf{Abstract} (line 65): Changed to ``CRBN$\to$VHL E3-unseen transfer.''
  \change{Line 65.}

  \item \textbf{Introduction} (line 136): ``VHL-targeted proteins fail below random (0.396 AUROC)---both limitations prominently acknowledged.''
  \change{Line 136.}

  \item \textbf{Section heading} (line 635): ``E3 generalization is limited to major ligases.''
  \change{Line 635.}

  \item \textbf{E3-unseen vs target-unseen paradox explained} (lines 636--638): CRBN$\to$VHL E3-unseen (0.811) succeeds because it tests E3-level transfer; VHL target-unseen (0.396) fails because it tests protein-level generalization---different axes.
  \change{Lines 636--638.}

  \item \textbf{VHL root cause analysis} (lines 640--648): Ruled out class imbalance (36.9\% vs 38.4\%) and dataset size (91 vs 125 proteins); identified structural heterogeneity as the likely cause. Explicit warning: ``practitioners should expect substantially lower performance than the overall 0.646 mean when deploying on VHL-targeted proteins, and near-random performance for rare E3s.''
  \change{Lines 640--648.}

  \item \textbf{Per-E3 Table~5} (lines 651--666): Quantitative breakdown showing CRBN 0.758, VHL 0.396, cIAP1 0.417.
  \change{Table~5 at line 654.}

  \item \textbf{Per-E3 Figure~5} (lines 668--674): Visual comparison of per-E3 performance.
  \change{Figure~5 at line 673.}

  \item \textbf{Conclusion} (line 783): ``CRBN$\to$VHL transfer (0.811 AUROC).''
  \change{Line 783.}
\end{itemize}

\bigskip

%% ====================================================================
%%  M5
%% ====================================================================

\reviewerbox{M5}{Validation setup should better align with the target-unseen setting, or explain why the current setup is sufficient.}

\resp We now explicitly acknowledge this as a methodological limitation in the Training Procedure section (\S3.5):

\begin{quote}\itshape
``Validation-split AUROC does not reliably predict target-unseen test AUROC\ldots\ Ideally, we would use a target-unseen validation split for hyperparameter selection; however, this would further reduce the already-limited training data (155 proteins $\to$ ${\sim}$120 training, ${\sim}$15 validation, ${\sim}$20 test).''
\end{quote}

We state that the chosen hyperparameters (LR=$5\times10^{-4}$, dropout=0.05) may not be optimal for target-unseen generalization, and that this likely contributes to the observed variance.

\change{Lines 326--330 (``Validation-test mismatch'' paragraph in \S3.5 Training); Line 323 (hyperparameter values noted); Limitation item~4 at line 769 (``Validation-test mismatch'').}

\bigskip

%% ====================================================================
%%  M6 (Minor 1)
%% ====================================================================

\reviewerbox{Minor 1}{Dataset limitations should be stated more directly in main results.}

\resp A ``Dataset characteristics and limitations'' paragraph now appears at the very start of the Results section (\S4), \emph{before} any performance numbers are presented. It states:

\begin{itemize}[topsep=2pt,itemsep=1pt]
  \item Only 155 unique protein targets (effective diversity for target-unseen).
  \item E3 distribution: CRBN 62\%, VHL 35\%, remaining 8 E3s combined 3\%.
  \item These constraints fundamentally bound generalization performance.
\end{itemize}

\change{Lines 431--436 (``Dataset characteristics and limitations'' paragraph); also Dataset statistics Table~2 at line 353 (``dominated by two E3 ligases, reflecting the current state of PROTAC research'').}

\bigskip

%% ====================================================================
%%  M7 (Minor 2)
%% ====================================================================

\reviewerbox{Minor 2}{The lack of external validation is an important limitation.}

\resp Now stated as Limitation item~6:

\begin{quote}\itshape
``No external validation: PROTAC-8K is the only publicly available benchmark with sufficient scale; generalization to data from other assay platforms or research groups is unvalidated.''
\end{quote}

We also note in the introduction (line 136) that limitations are ``prominently acknowledged,'' and in the conclusion (lines 783--784) that results are ``not statistically superior to gradient boosting.''

\change{Limitation item~6 at line 771 (``No external validation''); Introduction line 136; Conclusion lines 783--784.}

\bigskip
\hrule
\bigskip

\section*{Summary of All Changes Relevant to Reviewer Y2rf}

\begin{center}
\small
\begin{tabular}{@{}p{3.5cm}p{8cm}@{}}
\toprule
\textbf{Concern} & \textbf{Revision (lines in \texttt{paper.tex})} \\
\midrule
Best seed over-emphasized (M1) & Abstract L65; Introduction L135--136; Table~1 L441--458; \S4.2 L460--462; Conclusion L783--785 \\
Target-unseen unstable (M2) & \S4.3 L615--625 (seed sensitivity + 6-seed); Table~1 L452; Figure~4 L632; $p$-values L461, L529, L621 \\
Lysine claim unsupported (M3) & \S5 L751--758; Figure~7 L762; Limitation~7 L772 \\
E3 claims too broad (M4) & Abstract L65; \S4.3 L635--648; Table~5 L654; Figure~5 L673; Conclusion L783 \\
Validation misaligned (M5) & \S3.5 L326--330; L323; Limitation~4 L769 \\
Dataset limits buried (M6) & \S4 L431--436; Table~2 L353 \\
No external validation (M7) & Limitation~6 L771; Introduction L136; Conclusion L783--784 \\
\midrule
\multicolumn{2}{@{}l}{\textit{New analyses added in revision:}} \\
\midrule
6-seed extended validation & Table~1 L452; \S4.3 L619--625; Figure~4 L632 \\
Calibration analysis & \S4.6 L705--730; Table~6 L714 \\
Per-E3 breakdown & Table~5 L654; Figure~5 L673 \\
VHL root cause analysis & \S4.3 L640--648 \\
BRD4 case study & \S4.5 L686--703 \\
Lysine paradox figure & Figure~7 L762 \\
Statistical significance & L461, L529, L621 \\
Honest performance assessment & \S5 L744--748 \\
\bottomrule
\end{tabular}
\end{center}

\end{document}
